\documentclass[11pt]{article}
\usepackage[margin=1in]{geometry}
\usepackage[T1]{fontenc}
\usepackage{lmodern}
\usepackage{microtype}
\usepackage{amsmath,amssymb,amsfonts}
\usepackage{booktabs}
\usepackage{array}
\usepackage{enumitem}
\usepackage{graphicx}
\usepackage{placeins}
\usepackage{xcolor}
\usepackage{hyperref}
\usepackage[numbers,sort&compress]{natbib}

\hypersetup{
  colorlinks=true,
  linkcolor=blue!55!black,
  citecolor=blue!55!black,
  urlcolor=blue!55!black
}

\setlist[itemize]{leftmargin=1.6em}
\setlist[enumerate]{leftmargin=1.8em}
\newcommand{\code}[1]{\texttt{#1}}
\newcommand{\R}{\mathbb{R}}
\newcommand{\E}{\mathcal{E}}
\newcommand{\tr}{\operatorname{tr}}
\newcommand{\diag}{\operatorname{diag}}
\newcommand{\vol}{\operatorname{vol}}
\newcommand{\argmin}{\operatorname*{arg\,min}}

\newenvironment{reviewbox}[1]
{\par\medskip\noindent\textbf{#1.}\ }
{\par\medskip}

\newcommand{\paperfigure}[4]{%
\begin{figure}[htbp]
\centering
\includegraphics[width=0.82\linewidth]{#1}
\caption{#2 Reproduced from #3, internal review only. #4}
\end{figure}
}


\title{Review: Gaussian Fields and Harmonic Functions for Semi-Supervised Learning}
\author{Writer-P05 polished review}
\date{2026-05-15}

\begin{document}
\maketitle

\begin{abstract}
Zhu, Ghahramani, and Lafferty's ICML 2003 paper
\emph{Semi-Supervised Learning Using Gaussian Fields and Harmonic Functions}
is a compact canonical statement of graph-based semi-supervised learning as a
Dirichlet problem on a weighted graph \citep{zhu2003semisupervised}. Labeled
examples are clamped boundary nodes, unlabeled examples are interior nodes, and
predictions are obtained by minimizing a quadratic graph energy. The resulting
function is harmonic: at every unlabeled vertex it equals a weighted average of
neighboring values. This review treats the paper as a source for finite-graph
harmonic extension, conductance-weighted interpolation, class-mass
normalization, and entropy-based weight learning. For SIMODS and \code{gflow},
the paper is most useful not as a direct model of current overlap-density
smoothing, but as a clean reference for how the choice of graph weights or
kernel conductances controls a Laplacian solve.
\end{abstract}

\section{Why This Paper Matters}

The paper is important because it turns the informal semi-supervised premise
that ``nearby points should have similar labels'' into a precise finite-graph
calculus. Given labeled points, unlabeled points, and a symmetric nonnegative
weight matrix \(W=(w_{ij})\), the learning problem is not first to estimate a
parametric decision boundary. It is to find the smoothest function on the graph
that agrees with the known labels. The labeled nodes are fixed boundary values;
the unlabeled nodes receive the interior values of a harmonic extension.

This is a particularly useful formulation for a conductance and kernel
Laplacian review because the method exposes the role of the edge weights with
unusual clarity. In the graph-energy view, large \(w_{ij}\) makes disagreement
between \(f_i\) and \(f_j\) expensive. In the Markov view, the row-normalized
matrix \(P=D^{-1}W\) controls random-walk hitting probabilities. In the
electrical-network view used explicitly by the authors, the same weight is a
conductance, so the solution is a voltage field with labeled classes acting as
terminals. These interpretations are mathematically equivalent for the
Dirichlet problem, but they emphasize different reuse patterns for later
systems.

The paper is also valuable because it is sober about the graph. Its authors do
not present the Gaussian affinity as universally correct; they repeatedly note
that the right graph may be unclear, and they devote a section to learning
feature-wise length scales. That makes the paper an early, concrete reminder
that graph construction is part of the model, not disposable preprocessing.

\section{Problem Setting and Reader Orientation}

The setup is finite and direct. There are \(l\) labeled examples and \(u\)
unlabeled examples, with \(n=l+u\) graph vertices. The paper begins with binary
labels \(y_i\in\{0,1\}\), orders labeled vertices first, and assumes a
connected weighted graph. A typical affinity for Euclidean data is
\begin{equation}
  w_{ij}
  =
  \exp\left\{
    -\sum_{d=1}^{m} \frac{(x_{id}-x_{jd})^2}{\sigma_d^2}
  \right\}.
  \label{eq:p05-gaussian-weight}
\end{equation}
The parameters \(\sigma_d\) are global feature-wise length scales: small
\(\sigma_d\) makes the graph weight sensitive to coordinate \(d\), while large
\(\sigma_d\) suppresses that coordinate. For text, the paper instead uses a
sparse 10-nearest-neighbor graph under cosine similarity and an exponential
cosine-distance weight. Thus the common abstraction is not a single kernel
formula, but a weighted graph whose entries encode the intended notion of
similarity.

The main unknown is a real-valued function \(f\) on vertices. For binary
classification, a labeled zero node is clamped to \(f_i=0\), and a labeled one
node is clamped to \(f_i=1\). Unlabeled values are allowed to lie between zero
and one. The smoothness objective is the quadratic graph energy
\begin{equation}
  E(f)
  =
  \frac{1}{2}\sum_{i,j} w_{ij}(f_i-f_j)^2.
  \label{eq:p05-energy}
\end{equation}
This is the key object for a non-specialist reader. It says that changing the
function across a strong edge is costly, while changing it across a weak edge is
comparatively cheap. If \(W\) is read as an electrical conductance matrix, the
same expression is energy dissipation up to the usual convention for counting
undirected edges.

The paper derives this objective from a Gaussian random field whose density is
proportional to \(\exp[-\beta E(f)]\), constrained to match the labels. Because
the energy is quadratic, the most probable configuration is unique under the
Dirichlet boundary condition, and it is also the mean of the field. The authors
focus mainly on this mean/mode rather than on posterior uncertainty, but the
probabilistic framing explains why a continuous relaxation of labels yields a
closed-form linear-algebraic solution where discrete multi-label random fields
often require hard combinatorial optimization.

\section{The Harmonic Construction}

Let \(\Delta=D-W\) be the positive combinatorial graph Laplacian, where
\(D_{ii}=d_i=\sum_j w_{ij}\). Minimizing \eqref{eq:p05-energy} subject to fixed
labeled values gives the harmonic equation on unlabeled nodes:
\begin{equation}
  \Delta f = 0 \quad \text{on } U,
  \qquad
  f_L=y_L.
\end{equation}
Equivalently, each unlabeled value is the weighted average of its neighbors:
\begin{equation}
  f_j=\frac{1}{d_j}\sum_i w_{ij}f_i,
  \qquad j\in U.
  \label{eq:p05-harmonic-average}
\end{equation}
This local formula is the simplest way to understand harmonic extension. The
unknown value at a node is pulled toward the boundary labels through every path
in the graph, with each pull mediated by the conductances or affinities along
the way.

Writing labeled vertices before unlabeled vertices and partitioning
\[
W=
\begin{pmatrix}
W_{LL} & W_{LU}\\
W_{UL} & W_{UU}
\end{pmatrix},
\]
the paper's closed-form solution is
\begin{equation}
  f_U = (D_{UU}-W_{UU})^{-1}W_{UL}f_L.
  \label{eq:p05-dirichlet-solve}
\end{equation}
With \(P=D^{-1}W\), the same solution can be written as
\[
  f_U=(I-P_{UU})^{-1}P_{UL}f_L.
\]
The first expression foregrounds the unnormalized Dirichlet Laplacian solve; the
second foregrounds an absorbing Markov chain. Both are in the paper. Neither is
the present overlap-density smoother in \code{gflow}; for SIMODS they should be
read as finite-graph analogues and comparators, not as a description of current
implementation semantics.

\paperfigure{../../paper_figures/P05_F02_page03.png}
{Source Figure 2 crop shows the method's most compact visual argument:
with only a few large labeled points, harmonic energy minimization follows the
graph geometry of bands and spirals rather than a Euclidean nearest-neighbor
shortcut.}
{Zhu, Ghahramani, and Lafferty (2003)}
{The figure is useful here because it makes the boundary-value character of the
method visible: labels propagate through the weighted graph, not through an
explicit separating surface.}

\FloatBarrier

\section{Interpretations: Walks, Voltages, Heat, and Spectra}

One strength of the paper is that it does not leave the linear system as an
opaque algebraic recipe. It offers several equivalent readings.

In the random-walk view, an unlabeled node starts a walk with one-step
transition probabilities \(P_{ij}=w_{ij}/d_i\). Labeled nodes are absorbing
boundaries. Then \(f_i\) is the probability that a walk starting at \(i\) first
hits a positive labeled node before a negative labeled node. The harmonic
equation \eqref{eq:p05-harmonic-average} is exactly the first-step recursion for
that hitting probability.

In the electrical-network view, graph edges are resistors with conductance
\(W\). Nodes labeled one are connected to a positive voltage source; nodes
labeled zero are connected to ground. The unlabeled values \(f_U\) are the
resulting voltages, and the harmonic solution minimizes energy dissipation. This
is the paper's cleanest support for conductance language: outside this paragraph
it is safer to say ``weight'' or ``affinity, interpretable as conductance,''
because the graph is also used as a similarity and transition object.

In the heat-kernel and graph-kernel view, the paper uses the positive
combinatorial Laplacian with the heat kernel
\[
  K_t=e^{-t\Delta}.
\]
For the unlabeled subgraph with Dirichlet boundary conditions, it uses
\begin{equation}
  K'_t=e^{-t\Delta_{UU}},
  \qquad
  G=\int_0^\infty K'_t\,dt
    =\int_0^\infty e^{-t\Delta_{UU}}\,dt
    =(D_{UU}-W_{UU})^{-1}.
  \label{eq:p05-green}
\end{equation}
This Green's function is then the kernel behind the harmonic classifier:
\[
  f_U = G W_{UL}f_L.
\]
The audit correction is important here: the PDF's heat-kernel sign is negative
for the positive Laplacian. Using \(e^{t\Delta}\) would reverse the convention
and would not yield the inverse formula above.

The authors also compare the method to normalized cuts through the Rayleigh
quotient
\[
  R(f)=\frac{f^\top \Delta f}{f^\top D f}.
\]
The paper prints the numerator as a full double sum
\(\sum_{ij}w_{ij}(f_i-f_j)^2\). Under the standard symmetric-\(W\) convention,
\[
  f^\top \Delta f
  =
  \frac{1}{2}\sum_{i,j}w_{ij}(f_i-f_j)^2.
\]
The missing factor of two is immaterial for the Rayleigh quotient and the
generalized eigenproblem, but it matters if formulas are reused verbatim in a
conductance review where edge-counting conventions must stay consistent with
the energy in \eqref{eq:p05-energy}.

\section{Class Priors and Class-Mass Normalization}

The bare harmonic solution produces soft values, but the obvious decision rule
is to threshold at \(1/2\). The paper shows why that can be a poor final step.
If the estimated graph geometry is biased, the harmonic values may all be
shifted toward one class even when their relative ordering still carries useful
information. The authors' digit experiments are the striking case: for ``1'' vs.
``2'', many unlabeled values lie close to one, so thresholding labels too many
points as digit one.

Class-mass normalization (CMN) is their simple correction. If the desired prior
class proportions are \(q\) for class one and \(1-q\) for class zero, CMN compares
the prior-weighted and mass-normalized scores
\[
  \frac{q f_i}{\sum_{j\in U}f_j}
  \quad \text{and} \quad
  \frac{(1-q)(1-f_i)}{\sum_{j\in U}(1-f_j)}.
\]
The node is assigned to class one when the first score is larger. This does not
change the harmonic extension; it changes the decision rule after the graph has
propagated label information. The distinction is useful for SIMODS because it
separates an operator question from a calibration or class-balance question.

\paperfigure{../../paper_figures/P05_F03_F04_page06.png}
{Crop containing source Figures 3 and 4 summarizes the paper's
empirical case for CMN and graph-based propagation. The top row shows digit
tasks, including a failure of plain thresholding and a gain from combining a
voted perceptron with graph harmonic minimization. The lower row shows
20-newsgroups document tasks on sparse 10-nearest-neighbor cosine graphs.}
{Zhu, Ghahramani, and Lafferty (2003)}
{This combined crop is included because it links two review themes:
class-prior correction is crucial for some graph geometries, while the same
finite-graph machinery transfers beyond pixel-space Gaussian weights.}

\FloatBarrier

\section{External Classifiers and Softened Evidence}

The paper's graph-only solution assumes the manifold structure inferred from
features is sufficient. Section 5 acknowledges that this is often too austere.
An external classifier can be incorporated by attaching a labeled ``dongle''
node to every unlabeled node. If the external classifier predicts \(h_U\) and
the probability of jumping from an unlabeled node to its dongle is \(\eta\), the
paper gives
\[
  f_U
  =
  \left(I-(1-\eta)P_{UU}\right)^{-1}
  \left((1-\eta)P_{UL}f_L+\eta h_U\right).
\]
This formula is worth preserving because it shows how vertex-level evidence can
be mixed with graph propagation without abandoning the harmonic framework. The
authors also note that if labeled data are noisy, one could attach dongles to
labeled nodes as well and move labels onto those auxiliary nodes. That remark
marks a limitation of the main method: its default boundary values are hard,
not noisy observations.

\section{Entropy and Weight Learning}

The graph weights are the method's main modeling choice. Section 6 proposes to
learn the Gaussian length scales in \eqref{eq:p05-gaussian-weight} by minimizing
average entropy over unlabeled predictions:
\[
  H(f)=\frac{1}{u}\sum_{i\in U}
  \left[-f_i\log f_i-(1-f_i)\log(1-f_i)\right].
\]
The rationale is heuristic but coherent: a good graph, constrained by the known
labels, should yield confident interior values close to zero or one. The
authors differentiate the harmonic solution with respect to the length scales
and optimize by gradient descent.

The paper also exposes the danger of this objective. Without smoothing, entropy
has a degenerate minimum as \(\sigma\to 0\). The graph becomes too local, and
predictions collapse toward a hard nearest-neighbor-like propagation procedure.
To avoid this nuisance minimum, the authors replace the Markov transition matrix
by a PageRank-style smoothed matrix
\[
  \widetilde{P}=\epsilon U+(1-\epsilon)P,
\]
where \(U\) is the uniform transition matrix. The smoothing is not a decorative
technicality; it is what makes entropy minimization a usable weight-learning
criterion in their toy example.

\paperfigure{../../paper_figures/P05_F05_page07.png}
{Paper Figure 5 crop isolates the weight-learning issue. With no
smoothing, entropy falls toward the unhelpful \(\sigma\to 0\) limit and the
tighter grid invades the sparser one. With \(\epsilon=0.01\), the useful
minimum occurs near \(\sigma=0.67\), and the two-grid structure is recovered.}
{Zhu, Ghahramani, and Lafferty (2003)}
{This is the paper's clearest warning for kernel-conductance selection: making
a kernel very local can optimize a formal confidence criterion while degrading
the intended interpolation geometry.}

\FloatBarrier

On the digit task, the learned feature-wise \(\sigma_i\) values act as feature
selection. Small \(\sigma_i\) makes a pixel dimension influential; large
\(\sigma_i\) suppresses it. The paper reports that entropy falls from 0.6931 to
0.6542 bits after learning, with CMN accuracy increasing from
\(97.25\pm0.73\%\) to \(98.56\pm0.43\%\), and threshold accuracy increasing
from \(94.70\pm1.19\%\) to \(98.02\pm0.39\%\).

\paperfigure{../../paper_figures/P05_F06_page08.png}
{Crop containing Paper Table 1 and Paper Figure 6 shows the learned
pixel-level length scales for the ``1'' vs. ``2'' digit task. The learned
parameters suppress some dimensions and emphasize others, compensating for
differences in within-class tightness.}
{Zhu, Ghahramani, and Lafferty (2003)}
{For this review, the useful point is not the specific digit pattern but the
modeling principle: length scales alter the graph itself, and therefore alter
the harmonic solution.}

\FloatBarrier

\section{Limitations and Transfer Risks}

Several limits are worth keeping explicit. First, the main derivation is a
finite-graph Dirichlet interpolation method, not a continuum convergence theory
for graph Laplacians on manifolds. The paper uses manifold language and spectral
connections, but it does not prove that \(\Delta\), \(P\), or \(G\) converge to a
particular differential operator.

Second, the displayed formulas are mostly binary. The authors state that CMN and
the matrix solution extend naturally to multi-label classification, and they
contrast efficient Gaussian-field inference with NP-hard discrete multi-label
random fields. Still, a synthesis should not import detailed multi-class
algebra that is not in the paper.

Third, the graph is both the strength and the vulnerability of the method. Poor
weights can produce badly shifted soft predictions, and the entropy-learning
criterion itself can become degenerate without smoothing. The paper's remedy is
practical and illuminating, but it remains a heuristic rather than a fully
principled model-selection theory.

Fourth, the empirical benchmarks are early-2000s scale: synthetic structures,
downsampled handwritten digits, and 20-newsgroups text tasks. They demonstrate
the mechanics and promise of harmonic graph propagation, especially under scarce
labels, but they do not settle modern questions about calibration, robustness,
scalability, or graph construction in high-noise settings.

\section{Implications for SIMODS and \code{gflow}}

For SIMODS and \code{gflow}, the reusable core is the Dirichlet map
\[
  (W,f_L)\mapsto f_U.
\]
Once \(W\) is chosen, the solution follows from a Laplacian linear system or,
equivalently, from absorbing random walks. Changing the conductance or kernel
choice changes the interpolation even when graph topology and boundary values
are held fixed. That is precisely the conceptual pressure point for comparing
overlap-density conductance, inverse-length conductance, and kernel-based
conductance variants.

The paper should not be cited as evidence for current
\code{fit.rdgraph.regression()} overlap-density smoothing. Zhu et al. use
Gaussian feature-space weights for digits and exponential cosine-distance
weights on sparse text graphs. The transfer to SIMODS is therefore structural:
weighted graph regression by harmonic extension, not formula identity. In
particular, the paper supports language such as ``weights interpreted as
conductances in a Dirichlet energy'' and ``kernel choice controls label
propagation,'' while it does not support claims about overlap-density or
Riemannian-complex conductance.

The most useful design lesson is that graph construction, normalization, and
post-solve calibration are separate decisions. The unnormalized Laplacian solve
\((D_{UU}-W_{UU})^{-1}W_{UL}f_L\) and the row-normalized Markov form
\((I-P_{UU})^{-1}P_{UL}f_L\) are equivalent views of the same paper method, but
they foreground different diagnostics. The energy view asks where sharp
function changes are expensive; the Markov view asks how mass flows to labeled
boundaries; the CMN layer asks whether the resulting soft masses match known
class proportions. SIMODS comparisons should keep those layers distinct.

\section{What To Reuse}

The strongest reusable elements are the following. Use the graph energy
\eqref{eq:p05-energy} as a clean statement of conductance-weighted smoothness.
Use the harmonic average \eqref{eq:p05-harmonic-average} as the local
interpretation of interpolation. Use the Dirichlet solve
\eqref{eq:p05-dirichlet-solve} and Green's function \eqref{eq:p05-green} as the
finite-graph operator formulas. Use CMN as an example of post-hoc mass
calibration that does not alter the graph solve. Use Figure 5's smoothing story
as a warning that confidence objectives can reward pathological overlocal
kernels unless the transition process is regularized.

For wording, the most faithful formulation is: Zhu et al. define graph
weights/affinities that, in the electrical interpretation, are conductances.
That phrasing preserves the paper's terminology while still making the
conductance implication available for SIMODS.

\section*{Provenance}

Primary source: the local canonical P05 PDF in
\code{sources/pdf/}, read in full, including the figure/table crops corresponding to
Figures 1--6. Archival scaffolding: the local P05 Markdown memo in
\code{paper\_memos/} and the P05 audit memo in \code{audits/}. Figure
screenshots are internal-review captures listed in
\code{paper\_figure\_screenshots.yml} and included from
\code{../../paper\_figures/}. The original three-node SVG was reviewed as a
derived explanatory aid but not included here, since this LaTeX template does
not configure SVG inclusion.

\bibliographystyle{plainnat}
\bibliography{../../references}

\end{document}
